\section{Agent behavior: complete transition audit and selected examples}
\label{app:agent-behavior}

We separate a descriptive audit of every common-horizon episode from selected
examples of decision-making. The evidence is the 360 recorded handoffs in all
90 episodes through query 4, joined by geometry hash to their CFD grades.
\texttt{data/agent\_handoffs.csv} contains these decision rationales and the
additional steps used below; it does not contain full model reasoning streams.
A rationale describes the agent's stated basis for a decision, not an
independently validated physical explanation.

\subsection{Outcomes following an unsuccessful evaluation}

Every four-query episode has three adjacent query pairs. Of the 270 pairs,
64 begin with a request that returned no complete score. Table~\ref{tab:agent-transitions}
counts what happens immediately afterward. A changed geometry is identified
by its content hash; it does not, by itself, demonstrate that feedback caused
the change. Query-4 failures have no successor within the common horizon.

\begin{table}[htbp]
\centering\small
\caption{All transitions following a failed evaluation within queries 1--4.
Counts condition on different preceding failures and are correlated within
episodes; they are not a causal comparison of recovery strategies.}
\label{tab:agent-transitions}
\begin{tabular}{lrrrr}
\toprule
Method & Transitions & Episodes & New geometry & Next valid \\
\midrule
Open-ended research & 27 & 21 & 27 & 19 \\
Planning & 20 & 14 & 18 & 10 \\
Reflection & 17 & 15 & 17 & 15 \\
\bottomrule
\end{tabular}
\end{table}

This shows that the traces contain both successful and unsuccessful next
steps, rather than only successful anecdotes. It cannot attribute recovery
to a particular heuristic or prompt. The accompanying script emits every
transition so that alternative descriptive summaries can be checked.

\subsection{Selected decisions and their observed outcomes}

\textbf{A076, open-ended, low load: constrain extrapolation, then isolate a
change.} Query 1 tested a 47.73-degree sweep wing with predicted $L/D$ 20.64
and failed the mesh gate. Query 2 reduced sweep to 37.10 degrees and used a
dataset-trained mesh-quality estimate; CFD returned 16.04 against a surrogate
prediction of 18.30. Query 3 had a higher surrogate prediction, 19.08, but a
lower CFD score, 15.74. The next steps demonstrate both a failed adjustment
and a more controlled test (Table~\ref{tab:agent-case-a076}).

\begin{table}[htbp]
\centering\small
\caption{A076's subsequent decisions. Queries 5--8 belong to the selected
open-ended extension, not the common four-query comparison.}
\label{tab:agent-case-a076}
\begin{tabular}{cp{0.68\linewidth}r}
\toprule
Query & Recorded change & CFD $L/D$ \\
\midrule
4 & Increase q2 sweep to 38.3 degrees and reduce root thickness & timeout \\
5 & Near-repeat of q4, with root adjustment changed by 0.002 & timeout \\
6 & Return to q2 and change only sweep to 38.3 degrees & 16.35 \\
7 & Add a smaller root-thickness change & 16.64 \\
8 & Increase sweep to 39.0 degrees & 16.77 \\
\bottomrule
\end{tabular}
\end{table}

The q6 rationale explicitly tests whether sweep or thickness contributed to
the timeout. This is an example of using a controlled change after failure,
not proof of the timeout's physical cause or of the heuristic's general value.

\textbf{A080, open-ended, nominal load: revise a numerical optimizer's
objective.} The agent used CMA-ES on AeroTransformer. Its first two predictions
were 32.76 and 35.70, while CFD returned 32.06 and 33.15. It then used a
surrogate-bias correction based on distance to the source data: the next two
CFD scores were 33.4018 and 33.4014. The q4 rationale specifies the penalized
objective $y-5.083d$, where $d$ is nearest-source distance in the agent's
normalized coordinates. This illustrates optimizer orchestration and a
feedback-dependent change, but fitting a correction from two observations is
weak evidence for calibration, and the fourth step did not improve the score.

\textbf{A006, planning, high load: add a geometric screen.} The first wing
scored 28.52. A local refinement and then a smaller-radius refinement both
failed meshing. The executor subsequently computed a local surface-triangle
quality measure and screened candidates at the first wing's quality level.
Its fourth wing scored 28.55. The tool trace records the screening computation,
and the handoff names the threshold (0.0114), the candidate value (0.01232),
and the earlier failed candidates' values (at most 0.0106). The gain over q1
is only 0.03. Moreover, this episode's handoffs explicitly report substituting
an executable procedure for placeholder plans; it demonstrates executor
adaptation, not successful planning in isolation.

\textbf{A054, reflection, nominal load: distinguish predictors.} Four valid
scores increased from 28.79 to 31.14. On q2--q4, a supplementary predictor
reported 34.11, 34.35 and 34.84, about 12\% above CFD, while AeroTransformer
stayed within 0.5\%. The later rationales label the local predictor diagnostic
only. Building another model need not provide better candidate selection;
its disagreement should itself be tested against CFD.

\textbf{A035, reflection, high load: test predictor disagreement.} The q2
handoff prefers a candidate from the 2110 family over the higher-surrogate
5590 family because the latter conflicts with the parent's observed
high-Mach performance. At q4, the agent selects a 5590-family candidate as
an explicit test of disagreement between AeroTransformer (24.85) and its
local ensemble (22.91). CFD returns 24.89. The sequence shows how additional
modeling can discourage a productive direction, and how a subsequent
experiment can challenge that judgment; it does not isolate the ensemble's
effect from the other evidence used in selection.

\subsection{Implementation fidelity and limits of the lessons}

Twenty of the 120 planning handoff rationales, spanning seven episodes,
explicitly contain the word ``placeholder.'' This reproducible lexical count
is a lower bound on explicit mentions, not a classification of all plans.
Together with A006, it shows that valid structured output does not ensure an
actionable plan. Execution can supply substantive reasoning omitted by planning.

The campaign also used bounded schema or finalization recoveries. By query 4,
there were ten successful recovery decisions in nine episodes (six planning,
three reflection); the full extension adds two decisions in two open-ended
episodes. A006 and A054 each had a q1 recovery. All episodes remain included
in the reported outcomes. \texttt{data/agent\_recoveries.csv} identifies these
events; recovery assistance, additional stages and differing model usage are
part of the tested implementations, not effects isolated by randomization.

The actionable hypotheses are to preserve a verified fallback, make failure
diagnostics change candidate generation or screening, test surrogate
disagreement with observed CFD, and validate plans for executability. These
are motivated by the records and baseline comparisons. Demonstrating that any
one improves agent performance requires a prospective ablation; this audit
does not supply that causal evidence.
